\section{Statements become events}

Let \(\Omega\) be fixed. A statement about the situation becomes mathematical when we can say exactly for which elementary outcomes it is true. The set of all such outcomes is called an event.

\begin{definitionbox}
An \textbf{event} is a subset of the sample space \(\Omega\). If \(A\subseteq\Omega\), then \(A\) consists of precisely those elementary outcomes for which a certain statement is true.
\end{definitionbox}

This definition may look simple, but it is a major conceptual move. It tells us that statements and sets can play the same role once \(\Omega\) has been fixed:
\[
\text{statement about the situation}
\quad\longleftrightarrow\quad
\text{subset of }\Omega.
\]
The statement gives meaning. The subset gives structure.

\begin{examplebox}
Suppose two dice are rolled and order is recorded. Then
\[
\Omega=\{(i,j): i,j\in\{1,2,3,4,5,6\}\}.
\]
The statement
\[
\text{the sum is }7
\]
selects the event
\[
A=\{(1,6),(2,5),(3,4),(4,3),(5,2),(6,1)\}.
\]
The statement
\[
\text{the first die is greater than the second}
\]
selects the event
\[
B=\{(i,j)\in\Omega:i>j\}.
\]
The notation \(B=\{(i,j)\in\Omega:i>j\}\) is not a shortcut for avoiding thought. It records the rule by which membership in the event is decided.
\end{examplebox}

\section{Representing events as regions}

A set of outcomes can be written as a list, but lists are not always the best representation. When \(\Omega\) has a visible structure, an event should often be represented as a region inside that structure. This is not merely a visual aid. It is a way of seeing relations.

For two dice, the sample space is naturally represented as a \(6\times 6\) table:
\[
\begin{array}{c|cccccc}
 & 1&2&3&4&5&6\\
\hline
1&.&.&.&.&.&.\\
2&.&.&.&.&.&.\\
3&.&.&.&.&.&.\\
4&.&.&.&.&.&.\\
5&.&.&.&.&.&.\\
6&.&.&.&.&.&.
\end{array}
\]
Each cell represents one ordered pair \((i,j)\). The row records the first die, and the column records the second die.

The event \(B=\{(i,j):i>j\}\) occupies the cells below the diagonal:
\[
\begin{array}{c|cccccc}
 & 1&2&3&4&5&6\\
\hline
1&.&.&.&.&.&.\\
2&X&.&.&.&.&.\\
3&X&X&.&.&.&.\\
4&X&X&X&.&.&.\\
5&X&X&X&X&.&.\\
6&X&X&X&X&X&.
\end{array}
\]
The diagonal itself corresponds to \(i=j\). The region above the diagonal corresponds to \(i<j\). The geometry of the table reveals the structure of the statement.

\begin{remarkbox}
A representation is good when it preserves the distinctions relevant to the statement and makes relations visible. A poor representation may hide the structure that the problem is trying to teach.
\end{remarkbox}

\section{Reading a region as a statement}

The translation works in both directions. A verbal statement selects a set, but a marked set also asks to be interpreted as a statement.

\begin{examplebox}
In the two-dice table, suppose the marked cells are exactly the diagonal:
\[
\begin{array}{c|cccccc}
 & 1&2&3&4&5&6\\
\hline
1&X&.&.&.&.&.\\
2&.&X&.&.&.&.\\
3&.&.&X&.&.&.\\
4&.&.&.&X&.&.\\
5&.&.&.&.&X&.\\
6&.&.&.&.&.&X
\end{array}
\]
This region can be described in several equivalent ways:
\[
\text{the two dice show the same number},
\]
\[
\text{the first result equals the second result},
\]
\[
\{(i,j)\in\Omega:i=j\}.
\]
The ability to move between these descriptions is part of understanding the formalism.
\end{examplebox}

Students often think that formal notation is a translation from language into symbols. That is only half of the work. One must also be able to read symbols and diagrams back into language. A set without interpretation is mute; a sentence without a set is vague. Mathematics becomes powerful when both sides illuminate each other.
